\chapter{Resolution of Cosmological Anomalies: Dark Energy, Dark Matter, and the Hubble Tension}
\label{chap:cosmological_anomalies}


\section{The Hubble Artifact and the Photon Energy-Relaxation Length}
\label{sec:hubble_artifact}

When Edwin Hubble discovered the relationship between redshift and distance, standard cosmology interpreted the observed distance--redshift slope as evidence for the expansion of space. Tension Field Theory (TFT) assigns a different mechanical meaning to this observation. As established in \hyperref[chap:cosmic_redshift]{Chapter~\ref*{chap:cosmic_redshift}}, the measured low-redshift slope contains the combined projection of source--observer motion, gravitational endpoint loading, and coherent photon scale relaxation through the Active Vacuum.

The observational Hubble parameter therefore characterizes a measured spectral-shift gradient rather than a fundamental expansion rate. In the relaxation-dominated low-redshift reference branch, where systematic kinematic and gravitational gradients are negligible, the Chapter~\ref*{chap:cosmic_redshift} relation reduces to
\begin{equation}
    \Gamma_{\mathrm D}
    \simeq
    H_0^{(\mathrm{slope})}.
\end{equation}

Using the local observational reference value,
\begin{equation}
    H_0^{(\mathrm{obs})}
    \simeq
    73.04\ \mathrm{km\,s^{-1}\,Mpc^{-1}},
\end{equation}
this specialization defines the characteristic relaxation time
\begin{equation}
    t_{\mathrm{rel}}
    \equiv
    \frac{1}{\Gamma_{\mathrm D}}
    \simeq
    \frac{1}{H_0^{(\mathrm{obs})}}
    \approx
    13.39\ \mathrm{Gyr}.
\end{equation}

This timescale is the coherent-relaxation clock of the Active Vacuum. It is distinct from the thermodynamic age of the present baryonic phase, whose chronology begins with Catastrophic Decompression.

The corresponding coherent-relaxation length is the Photon Energy-Relaxation Length,
\begin{equation}
    R_{\mathrm U}
    \equiv
    \frac{c_{\mathrm E}}{\Gamma_{\mathrm D}}.
\end{equation}

Under the same low-redshift reference specialization,
\begin{equation}
    R_{\mathrm U}
    \simeq
    \frac{c_{\mathrm E}}{H_0^{(\mathrm{obs})}}
    \simeq
    1.2665\times10^{26}\ \mathrm{m}
    \simeq
    4104.5\ \mathrm{Mpc}.
\end{equation}

To understand the mechanical meaning of this distance, consider the coherent-action-preserving photon-relaxation law developed in
\hyperref[app:photon_dispersion_relaxation]{Appendix~\ref*{app:photon_dispersion_relaxation}}:
\begin{equation}
    E_{\mathrm{obs}}
    =
    E_{\mathrm{emit}}
    \exp\left(
        -\frac{\Gamma_{\mathrm D}x}{c_{\mathrm E}}
    \right)
\end{equation}
in the gravity-neutral, constant-rate limit.

At
\begin{equation}
    x
    =
    R_{\mathrm U}
    =
    \frac{c_{\mathrm E}}{\Gamma_{\mathrm D}},
\end{equation}
the exponent becomes unity,
\begin{equation}
    E_{\mathrm{obs}}
    =
    E_{\mathrm{emit}}e^{-1}
    \approx
    0.3679E_{\mathrm{emit}}.
\end{equation}

Thus $R_{\mathrm U}$ is the one-$e$-fold coherent-relaxation distance: after one such path length, approximately $63.2\%$ of the original coherent carrier energy has been transferred out of the surviving photon branch into dynamically accessible finite-energy receiving degrees of freedom.

\vspace{0.5em}

\phantomsection
\label{sec:kinematic_projection_artifact}

\noindent\textbf{The Kinematic Projection Artifact}

In standard FLRW cosmology, the Hubble radius $c/H_0$ is the distance at which the formal recession speed $H_0d$ reaches $c$. It is a kinematic scale of the expanding coordinate system rather than a material wall.

TFT assigns a different physical meaning to the corresponding macroscopic length. The scale $R_{\mathrm U}=c_{\mathrm E}/\Gamma_{\mathrm D}$ is not a recession-velocity threshold; it is the distance at which the coherent-relaxation exponent reaches unity.

A photon can therefore propagate through multiple $R_{\mathrm U}$ intervals while remaining on the coherent eigenbranch. The underlying Structural Reservoir extends vastly beyond this optical scale, and the finite thermodynamic history of the present matter--radiation cycle remains distinct from the pre-existing extent of the TFT continuum.


\section{The Accelerating Expansion Illusion: The Nonlinear Distance Law}

In 1998, observations of distant Type Ia supernovae revealed a nonlinear relation between redshift and apparent luminosity distance \cite{Riess_1998}. Standard cosmology represents this behavior through the time evolution of the metric scale factor and attributes the late-time acceleration to Dark Energy.

In TFT, the same observational problem is expressed through the combined mechanics of source motion, gravitational endpoint clocks, coherent photon scale relaxation, and photometric projection. The nonlinear distance law therefore arises from the mapping between the source state and the received radiation rather than from acceleration of the underlying Structural Reservoir.

For the gravity-neutral radial photometric branch derived in
\hyperref[chap:cosmic_redshift]{Chapter~\ref*{chap:cosmic_redshift}},
\begin{equation}
    1+z_{\mathrm{obs}}
    =
    (1+z_{\mathrm{kin}})
    (1+z_{\mathrm D}),
\end{equation}
and the luminosity distance is
\begin{equation}
    d_{\mathrm L}
    =
    r
    (1+z_{\mathrm{kin}})^2
    \sqrt{1+z_{\mathrm D}}.
\end{equation}

For constant coherent relaxation,
\begin{equation}
    \ln(1+z_{\mathrm D})
    =
    \frac{\Gamma_{\mathrm D}r}{c_{\mathrm E}},
\end{equation}
so that
\begin{equation}
    r
    =
    \frac{c_{\mathrm E}}{\Gamma_{\mathrm D}}
    \ln(1+z_{\mathrm D}).
\end{equation}

Substitution into the luminosity-distance relation gives
\begin{equation}
    d_{\mathrm L}
    =
    \frac{c_{\mathrm E}}{\Gamma_{\mathrm D}}
    (1+z_{\mathrm{kin}})^2
    \sqrt{1+z_{\mathrm D}}
    \ln(1+z_{\mathrm D}).
\end{equation}

Using
\begin{equation}
    1+z_{\mathrm D}
    =
    \frac{1+z_{\mathrm{obs}}}
    {1+z_{\mathrm{kin}}},
\end{equation}
the same relation may be written directly in terms of the total observed spectral shift:
\begin{equation}
    d_{\mathrm L}
    =
    \frac{c_{\mathrm E}}{\Gamma_{\mathrm D}}
    (1+z_{\mathrm{kin}})^{3/2}
    (1+z_{\mathrm{obs}})^{1/2}
    \ln\left(
        \frac{1+z_{\mathrm{obs}}}
        {1+z_{\mathrm{kin}}}
    \right).
\end{equation}

The logarithmic propagation coordinate therefore survives inside a broader nonlinear observational projection. Matter motion controls $z_{\mathrm{kin}}$, photon relaxation controls the accumulated logarithmic scale change, and gravitational endpoint loading enters through the full three-factor spectral law developed in Chapter~\ref*{chap:cosmic_redshift}.

The apparent acceleration of the supernova Hubble diagram is consequently interpreted in TFT as a nonlinear mapping between source dynamics and received radiation. The macroscopic matter-flow field, gravitational loading, and photon-relaxation history together determine the observable distance--redshift curve without requiring acceleration of the Structural Reservoir itself.


\section{The Hubble Tension: Acoustic Damping and the Standard Ruler}
\label{sec:hubble_tension_acoustic_damping}

A major tension in modern cosmology is the difference between the expansion rate measured from the local distance ladder,
\begin{equation}
    H_0
    \approx
    73.04\ \mathrm{km\,s^{-1}\,Mpc^{-1}},
\end{equation}
and the lower value inferred from the Cosmic Microwave Background within the $\Lambda$CDM acoustic calibration,
\begin{equation}
    H_0
    \approx
    67.4\ \mathrm{km\,s^{-1}\,Mpc^{-1}}.
\end{equation}

Standard cosmology infers the CMB value through the observed acoustic angular scale $\theta_*$ together with the calculated comoving photon--baryon sound horizon.

Within $\Lambda$CDM, the characteristic ruler is obtained from sound propagation through the early photon--baryon plasma and the evolving metric background, giving a comoving drag scale of approximately
\begin{equation}
    r_{\mathrm{drag}}
    \approx
    147\ \mathrm{Mpc}.
\end{equation}

TFT instead interprets the primordial acoustic structure within a globally static FCC-condensate medium and projects the observed angular pattern through the characteristic optical scale of that medium.


\subsection{Acoustic Propagation and Impedance in the Active Vacuum}
\label{sec:viscous_acoustic_damping}

In TFT, the primordial baryonic plasma evolved inside the restored State-I FCC Active Vacuum following Catastrophic Decompression. Its acoustic structure therefore developed within the coupled FCC--Condensate medium rather than within an empty background.

\begin{itemize}

    \item \textbf{The Waves:}
    Baryon Acoustic Oscillations are represented as longitudinal pressure structures propagating through the primordial baryonic plasma and the connected interstitial condensate.

    \item \textbf{The Impedance:}
    The Active Vacuum possesses the characteristic macroscopic impedance acceleration
    \begin{equation}
        a_{\mathrm D}
        =
        c_{\mathrm E}\Gamma_{\mathrm D}
        =
        \frac{c_{\mathrm E}^2}{R_{\mathrm U}}.
    \end{equation}
    This scale characterizes the FCC--Condensate response that also appears, through a distinct reactive projection, in galactic dynamics.

    \item \textbf{The Fossilization:}
    The acoustic pressure field evolves through this impedance-bearing medium until lattice annealing, recombination, and optical transparency freeze the primordial density pattern into the fossil matter distribution.

\end{itemize}

The acoustic ruler in TFT is therefore a physical structure of the State-I medium rather than a comoving coordinate length generated by metric expansion.


\subsection{The Euclidean Acoustic Resolution}
\label{sec:euclidean_acoustic_resolution}

The static TFT geometry assigns a physical Euclidean distance to the acoustic projection. The characteristic optical length is
\begin{equation}
    R_{\mathrm U}
    =
    \frac{c_{\mathrm E}}{\Gamma_{\mathrm D}}.
\end{equation}

Under the retained low-redshift reference specialization,
\begin{equation}
    \Gamma_{\mathrm D}
    \simeq
    H_0^{(\mathrm{obs})}
    \simeq
    73.04\ \mathrm{km\,s^{-1}\,Mpc^{-1}},
\end{equation}
giving
\begin{equation}
    R_{\mathrm U}
    \simeq
    13.39\ \mathrm{Gly}
    \simeq
    4104.498\ \mathrm{Mpc}.
\end{equation}

For the acoustic projection used in TFT, define
\begin{equation}
    D_{\mathrm{proj}}
    \equiv
    R_{\mathrm U}.
\end{equation}

Applying the observed CMB acoustic angular scale,
\begin{equation}
    \theta_*
    \simeq
    0.01041,
\end{equation}
gives the TFT Effective Acoustic Scale
\begin{equation}
    r_{\mathrm s}^{(\mathrm{TFT})}
    =
    D_{\mathrm{proj}}\theta_*
    =
    R_{\mathrm U}\theta_*.
\end{equation}

Numerically,
\begin{equation}
    r_{\mathrm s}^{(\mathrm{TFT})}
    =
    (4104.498\ \mathrm{Mpc})
    (0.01041)
    \approx
    42.73\ \mathrm{Mpc}.
\end{equation}

The static Euclidean projection therefore maps the observed acoustic angle onto a characteristic physical scale of approximately $42.73$ Mpc, substantially smaller than the $\Lambda$CDM comoving sound horizon of approximately $147$ Mpc. The difference follows directly from the distinct geometric distance assigned to the same observed angular structure.


\subsection{The Statistical Nature of the Effective Acoustic Scale}

The geometric projection
\begin{equation}
    r_{\mathrm s}^{(\mathrm{TFT})}
    =
    R_{\mathrm U}\theta_*
\end{equation}
defines the characteristic acoustic scale of the fossil rupture network rather than a universal diameter for every individual void or shell.

\hyperref[app:cosmic_energy_budget_and_chronology]{Appendix~\ref*{app:cosmic_energy_budget_and_chronology}} identifies the global rupture threshold as
\begin{equation}
    \beta_{\mathrm{rup}}
    =
    p_{\mathrm c}\beta_{\mathrm{IC}}
    \approx
    0.0516,
\end{equation}
corresponding to an aggregate live-boundary mass fraction
\begin{equation}
    \frac{M_{\mathrm p}}
    {M_{\mathrm{Res}}}
    \rightarrow
    0.0516
\end{equation}
at the onset of reservoir-spanning percolation.

The pre-rupture universe therefore contains a distributed population of State-II domains and saturated connecting bridges rather than one uniform spherical precursor.

The phase conversion obeys the mass-energy conservation relation
\begin{equation}
    V_{\mathrm{FCC}}\rho_{\mathrm{QG}}
    =
    V_{\mathrm{SC}}\rho_{\mathrm S}.
\end{equation}

Using the radial conversion factor derived in
\hyperref[sec:reservoir_radius]{Section~\ref*{sec:reservoir_radius}},
\begin{equation}
    R_{\mathrm{FCC}}
    \approx
    R_{\mathrm{SC}}
    \left(
        \frac{\rho_{\mathrm S}}
        {\rho_{\mathrm{QG}}}
    \right)^{1/3}
    \approx
    R_{\mathrm{SC}}
    \left(
        4.6098\times10^8
    \right),
\end{equation}
the compact precursor State-II domains map into vastly larger decompressed State-I regions.

Within this structural interpretation, the projected $42.73$ Mpc scale characterizes the dominant spatial frequency of the fossil rupture network. The observed range of void, wall, filament, and supercluster dimensions arises from the distribution of precursor domain sizes, overlap of neighboring decompression fronts, and subsequent gravitational evolution.


\section{Recalibrating the Kinematic Baseline}
\label{sec:recalibrating_kinematic_baseline}

The observed CMB acoustic angle acquires different physical scales depending on the geometric distance through which it is projected. The competing cosmological constructions may therefore be summarized as:

\begin{itemize}

    \item \textbf{Standard Model (Expanding Metric):}

    Damped Photon--Baryon Plasma

    $\to$

    Metric Comoving Distance

    ($D_{\mathrm M} \approx 14{,}000$ Mpc)

    $\to$

    Comoving Sound Horizon

    ($r_{\mathrm{drag}} \approx 147$ Mpc)

    $\to$

    CMB-Inferred Rate

    ($H_0 \approx 67.4$ km/s/Mpc).

    \item \textbf{TFT (Globally Static Dual Scale Dual-Phase Physical Vacuum):}

    Lattice-Coupled Acoustic Medium

    $\to$

    Euclidean Projection

    ($D_{\mathrm{proj}}\equiv R_{\mathrm U}\approx4104.5$ Mpc)

    $\to$

    TFT Effective Acoustic Scale

    ($r_{\mathrm s}^{(\mathrm{TFT})}\approx42.73$ Mpc)

    $\to$

    Local Hubble-Scale Projection

    ($H_0^{(\mathrm{obs})}\approx73.04$ km/s/Mpc).

\end{itemize}

Within TFT, the Hubble Tension is therefore interpreted as a calibration difference between a metric-comoving acoustic construction and the Euclidean projection of the same observed angular structure through the Active Vacuum. The resulting TFT Effective Acoustic Scale is
\begin{equation}
    r_{\mathrm s}^{(\mathrm{TFT})}
    \approx
    42.73\ \mathrm{Mpc}.
\end{equation}

\begin{figure}[H]

    \centering

    \includegraphics[width=0.9\textwidth]{Fig22_PartIII_CMB.pdf}

    \caption{\textbf{The Macroscopic Thermodynamic Projection and Acoustic Mapping.}
A Euclidean representation of Cosmic Microwave Background anisotropies within the TFT Active Vacuum.
\textbf{Photon Energy-Relaxation Scale:}
The characteristic one-$e$-fold propagation length is
$R_{\mathrm U}=c_{\mathrm E}/\Gamma_{\mathrm D}
\approx1.27\times10^{26}\ \mathrm{m}
\approx13.39\ \mathrm{Gly}$,
with the corresponding Photon Energy-Relaxation Time
$t_{\mathrm D}=\Gamma_{\mathrm D}^{-1}\approx13.39\ \mathrm{Gyr}$
in the low-redshift reference realization.
\textbf{Angular Projection:}
The observed acoustic angular scale
$\theta_*\approx0.01041$ rad
is projected across
$D_{\mathrm{proj}}\equiv R_{\mathrm U}$.
\textbf{Effective Acoustic Scale:}
This gives the TFT Effective Acoustic Scale
$r_{\mathrm s}^{(\mathrm{TFT})}
=R_{\mathrm U}\theta_*
\approx42.73$ Mpc,
compared with the $\Lambda$CDM comoving sound horizon of approximately $147$ Mpc.
\textbf{Relaxation-Rate Reference:}
In the low-redshift relaxation-dominated specialization,
$\Gamma_{\mathrm D}\simeq H_0^{(\mathrm{obs})}$,
with
$H_0^{(\mathrm{obs})}=73.04\ \mathrm{km\,s^{-1}\,Mpc^{-1}}$.
\textbf{Thermodynamic Impedance Signature:}
The observed $\pm30\,\mu\mathrm K$ structure is interpreted as the thermodynamic imprint of Fossil CMB radiation interacting with the rupture-generated density network and the receiving degrees of freedom of the Active Vacuum.}

    \label{fig:cmb_acoustic_damping}

\end{figure}


\section{The Thermodynamic Impedance Signature and the $\theta_*$ Anomaly}
\label{sec:thermodynamic_impedance_signature}

Standard cosmology interprets the characteristic
\begin{equation}
    100\theta_*
    \approx
    1.041
\end{equation}
angular scale measured by Planck as the fossil acoustic signature of the primordial photon--baryon plasma observed through the last-scattering geometry.

In TFT, the CMB is the fossil electromagnetic photon population produced by the rupture-and-annealing history developed in \hyperref[chap:nucleation_annealing]{Chapter \ref*{chap:nucleation_annealing}}. Its angular structure is therefore interpreted through the interaction of this fossil radiation with the large-scale density pattern and the finite-energy receiving response of the Active Vacuum.


\subsection{The Origin of the Fossilized Rings}

As derived in
\hyperref[sec:cosmic_web_formation]{Section~\ref*{sec:cosmic_web_formation}},
the pre-rupture State-II domains carry their inertial matter on live FCC/SC phase boundaries. When global confinement fails, rapid SC$\rightarrow$FCC decompression restores the State-I Active Vacuum and drives the phase-segregated baryonic matter outward with the expanding compressive and acoustic fronts.

The interiors of the decompressed domains become matter-depleted regions, while collisions and intersections of the outward-moving matter shells generate walls, filaments, and dense junctions. These fossilized density structures preserve the mechanical imprint of the Big Rupture.

Within the Euclidean acoustic projection above, the characteristic
\begin{equation}
    r_{\mathrm s}^{(\mathrm{TFT})}
    \approx
    42.73\ \mathrm{Mpc}
\end{equation}
scale represents the dominant projected spacing of this fossil rupture network and provides the TFT mechanical interpretation of the large-scale acoustic pattern.


\subsection{The Propagation-Relaxation Path Penalty}

The fossilized acoustic structures are physical density variations within the Active Vacuum. Radiation propagating through a dense historical shell encounters enhanced gravitational loading $\alpha_{\mathrm G}$, which modifies the local ray speed
\begin{equation}
    v_{\mathrm{ray}}
    =
    v_{\mathrm{ray}}
    \left[
        \omega,
        \alpha_{\mathrm G}
    \right].
\end{equation}

The propagation-relaxation depth is
\begin{equation}
    \ln(1+z_{\mathrm D})
    =
    \int_{\mathrm{path}}
    \frac{
        \Gamma_{\mathrm D}(\ell)
    }{
        v_{\mathrm{ray}}
        \left[
            \omega(\ell),
            \alpha_{\mathrm G}(\ell)
        \right]
    }
    \,d\ell.
\end{equation}

A reduction of the local ray speed therefore increases the elapsed propagation time and the accumulated coherent-relaxation depth across the same coordinate interval.

Dense fossil structures consequently alter the transmitted fossil-radiation field relative to propagation through matter-depleted void regions. At the same time, energy transferred from the coherent photon branch enters the finite-energy receiving sector and participates in local absorption--emission balance.

The observed thermal pattern is therefore generated by the combined spatial variation of transmitted fossil radiation and local receiver-mediated re-emission. Dense shell regions and matter-depleted void regions produce different thermodynamic histories, imprinting the positive and negative $\mu$K structure observed across the CMB sky.

\vspace{0.5em}

\noindent\textbf{Conclusion:}

In TFT, the CMB anisotropy field is a Thermodynamic Impedance Signature of the fossil rupture network. The angular pattern records the interaction between primordial electromagnetic radiation, the decompression-generated matter distribution, and the propagation-relaxation and finite-energy receiving response of the Active Vacuum.

\section{Kinematic Unification: Lattice Impedance ($a_{\mathrm D}$) and the Dark Sector}

A primary synthesis of Tension Field Theory is the connection between the structural Dark Sector and the macroscopic response of the Active Vacuum.

In \hyperref[part:one]{Part~\ref{part:one}}, the exact FCC geometry establishes the unperturbed structural partition
\begin{equation}
    \beta_{\mathrm{geo}}
    \approx
    0.74048,
    \qquad
    \beta_{\mathrm{IC}}
    \approx
    0.25952.
\end{equation}

Under the baryonic defect-conversion closure,
\begin{equation}
    \beta_{\Lambda,\mathrm{eff}}
    =
    \beta_{\mathrm{geo}}
    -
    \beta_{\mathrm b}.
\end{equation}

For
\begin{equation}
    \beta_{\mathrm b}
    \approx
    0.05,
\end{equation}
the remaining smooth nodal structural allocation is
\begin{equation}
    \beta_{\Lambda,\mathrm{eff}}
    \approx
    0.69048.
\end{equation}

The remaining nodal sector carries the smooth tensile background associated in TFT with the effective Dark-Energy contribution. The interstitial structural allocation occupies the connected condensate network, whose extended coherent response generates the galactic loading conventionally attributed to Dark Matter.

The two sectors are therefore structural components of the same State-I Active Vacuum, while their observable manifestations arise through distinct mechanical response channels.

The characteristic macroscopic impedance acceleration is
\begin{equation}
    a_{\mathrm D}
    =
    c_{\mathrm E}\Gamma_{\mathrm D}
    =
    \frac{c_{\mathrm E}^2}{R_{\mathrm U}}.
\end{equation}

Using the retained low-redshift reference scale,
\begin{equation}
    R_{\mathrm U}
    \simeq
    1.2665\times10^{26}\ \mathrm{m},
\end{equation}
gives
\begin{equation}
    a_{\mathrm D}
    \approx
    7.0964\times10^{-10}\ \mathrm{m\,s^{-2}}.
\end{equation}

The optical relaxation channel represents the dissipative scale-transfer projection of the Active Vacuum. Galactic dynamics probes a different, reactive response of the interstitial condensate. As developed in
\hyperref[app:cosmic_impedance]{Appendix~\ref*{app:cosmic_impedance}},
the corresponding galactic acceleration scale is
\begin{equation}
    a_{\mathrm{gal}}
    =
    \frac{a_{\mathrm D}}{2\pi}.
\end{equation}

Numerically,
\begin{equation}
    a_{\mathrm{gal}}
    \approx
    1.1294\times10^{-10}\ \mathrm{m\,s^{-2}}.
\end{equation}

This lies at the empirical acceleration scale
\begin{equation}
    a_0
    \approx
    1.1\text{--}1.2\times10^{-10}\ \mathrm{m\,s^{-2}}
\end{equation}
associated with the low-acceleration phenomenology of galactic rotation curves \cite{Milgrom_1983}.

The TFT dark-sector picture therefore has three linked mechanical layers: the residual FCC nodal allocation supplies the smooth tensile background, the interstitial condensate supplies the coherent reactive galactic response, and the same Active-Vacuum architecture supports the dissipative optical scale-transfer channel characterized by $\Gamma_{\mathrm D}$.

\vspace{0.5em}

\noindent\textbf{Conclusion:}

The apparent Dark-Energy and Dark-Matter phenomena are not independent substances in TFT. They are different structural and dynamical projections of the dual-phase Active Vacuum: the smooth FCC nodal sector carries the cosmological tensile background, while the interstitial condensate generates the extended galactic impedance response. Photon scale relaxation provides a separate dissipative projection of the same mechanical medium.


\vspace{0.5em}

\noindent\textbf{Global Scale Correspondence: Static Potential and Kinematic Relaxation}

At the Type-II saturation condition
\begin{equation}
    R_{\mathrm{QG}}
    =
    R_{\mathrm U},
\end{equation}
the corresponding saturation mass is
\begin{equation}
    M_{\mathrm{Sat}}
    =
    \frac{c_{\mathrm E}^2R_{\mathrm U}}{G}.
\end{equation}

Using
\begin{equation}
    R_{\mathrm U}
    =
    \frac{c_{\mathrm E}}{\Gamma_{\mathrm D}},
\end{equation}
this becomes
\begin{equation}
    M_{\mathrm{Sat}}
    =
    \frac{c_{\mathrm E}^3}
    {G\Gamma_{\mathrm D}}.
\end{equation}

The corresponding gravitational potential scale is
\begin{equation}
    \frac{
        GM_{\mathrm{Sat}}
    }{
        R_{\mathrm U}
    }
    =
    c_{\mathrm E}^2.
\end{equation}

The Active-Vacuum impedance relation gives independently
\begin{equation}
    a_{\mathrm D}
    =
    c_{\mathrm E}\Gamma_{\mathrm D},
    \qquad
    R_{\mathrm U}
    =
    \frac{c_{\mathrm E}}
    {\Gamma_{\mathrm D}},
\end{equation}
so
\begin{equation}
    a_{\mathrm D}R_{\mathrm U}
    =
    c_{\mathrm E}^2.
\end{equation}

The gravitational and relaxation scales therefore meet at the common kinematic capacity,
\begin{equation}
    \frac{
        GM_{\mathrm{Sat}}
    }{
        R_{\mathrm U}
    }
    =
    a_{\mathrm D}R_{\mathrm U}
    =
    c_{\mathrm E}^2.
\end{equation}

With the gravitational potential sign convention,
\begin{equation}
    -\frac{
        GM_{\mathrm{Sat}}
    }{
        R_{\mathrm U}
    }
    =
    -a_{\mathrm D}R_{\mathrm U}
    =
    -c_{\mathrm E}^2.
\end{equation}

This equality is the global TFT scale correspondence between Type-II gravitational saturation and one-$e$-fold optical relaxation. Static gravitational potential and dissipative photon transport remain distinct mechanical operations, but both terminate on the same invariant specific-energy scale $c_{\mathrm E}^2$.
